% GENERATED FILE. Edit canonical lessons, then run npm run generate:books.
% canonical-source-hash: fnv1a64:e427e3eb30cd8b65
% canonical-lessons: TA-C507-paniyitam, TA-C507-panipuri, TA-C507-ceyalalar, TA-C507-koppu, TA-C507-merparvaiyitu

\chapter{At the Office Desk}
\label{ch:ta-at-the-office-desk}
\hlchaptermodality{507}

\begin{chapteropening}
\textbf{By the end of this chapter:} \emph{I can talk about the place where I work.}

Complete the last lesson of chapter 507: I can talk about the place where I work.
\end{chapteropening}

\section[\texorpdfstring{\ta{பணியிடம்} (paṇiyiṭam)}{பணியிடம் (paṇiyiṭam)}]{\ta{பணியிடம்} (paṇiyiṭam) — a workplace}

\label{lesson:TA-C507-paniyitam}



\begin{quote}
Before the new one: say the Tamil for full-time, then the Tamil for an assistant.
\end{quote}

\subsection*{You'll want to know: \ta{பணியிடம்}}
\textbf{\ta{பணியிடம்}} — \emph{paṇiyiṭam} — \textquotedblleft{}a workplace\textquotedblright{}.

\begin{quote}
\ta{என்} \ta{பணியிடம்} \ta{வீட்டிலிருந்து} \ta{தூரமாக} \ta{இருக்கிறது}.
\end{quote}

\begin{quote}
\emph{eṉ paṇiyiṭam vīṭṭiliruntu tūramāka irukkiṟatu.}
\end{quote}

\begin{quote}
My workplace is far from home.
\end{quote}

Like any -\ta{அம்} noun it takes -\ta{அத்து}- before an ending: \ta{பணியிடத்தில்}, at the workplace.

\begin{cousinweb}
\textbf{\ta{பணியிடம்}} (\emph{paṇiyiṭam}) is a native compound of \ta{பணி} (\emph{paṇi}), \textquotedblleft{}work, service\textquotedblright{} (DEDR 3884), and \ta{இடம்} (\emph{iṭam}), \textquotedblleft{}place\textquotedblright{} (DEDR 434).
\end{cousinweb}

\begin{grammarlens}[title={What you've built}]
One more word for describing what you lived and felt.
\end{grammarlens}

\subsection*{Your turn}
\begin{itemize}
\raggedright
  \item \emph{Say it:} \emph{paṇiyiṭam}
  \item \emph{Say it:} \emph{paṇiyiṭam}, once more
  \item \emph{Say it:} \emph{eṉ paṇiyiṭam vīṭṭiliruntu tūramāka irukkiṟatu.} — My workplace is far from home.
  \item \emph{Recall:} say the Tamil for full-time, then the Tamil for an assistant, then say \emph{paṇiyiṭam} again
\end{itemize}

\begin{tcolorbox}[breakable,colback=teal!4,colframe=teal!35!black,title={Before you move on}]
What does \ta{பணியிடம்} mean? (\textquotedblleft{}A workplace\textquotedblright{}.) Say it once more.
\end{tcolorbox}
\section[\texorpdfstring{\ta{பணிபுரி} (paṇipuri)}{பணிபுரி (paṇipuri)}]{\ta{பணிபுரி} (paṇipuri) — to be employed; to work}

\label{lesson:TA-C507-panipuri}



\begin{quote}
Before the new one: say the Tamil for an assistant, then the Tamil for a workplace.
\end{quote}

\subsection*{You'll want to know: \ta{பணிபுரி}}
\textbf{\ta{பணிபுரி}} — \emph{paṇipuri} — \textquotedblleft{}to be employed; to work\textquotedblright{}.

\begin{quote}
\ta{என்} \ta{அப்பா} \ta{முப்பது} \ta{ஆண்டுகள்} \ta{ரயில்} \ta{நிலையத்தில்} \ta{பணிபுரிந்தார்}.
\end{quote}

\begin{quote}
\emph{eṉ appā muppatu āṇṭukaḷ rayil nilaiyattil paṇipurintār.}
\end{quote}

\begin{quote}
My father worked at the railway station for thirty years.
\end{quote}

Past (I …): \ta{பணிபுரிந்தேன்}, \emph{paṇipurintēṉ}.

A formal verb beside \ta{வேலை} \ta{செய்} (\emph{vēlai cey}). Where someone works takes -\ta{இல்}: \ta{வங்கியில்} \ta{பணிபுரிகிறார்}.

\begin{cousinweb}
\textbf{\ta{பணிபுரி}} (\emph{paṇipuri}) is a native compound: \ta{பணி} (\emph{paṇi}), \textquotedblleft{}work, service\textquotedblright{} (DEDR 3884), with \ta{புரி} (\emph{puri}), \textquotedblleft{}to do, to perform\textquotedblright{}, as the verb.
\end{cousinweb}

\begin{grammarlens}[title={What you've built}]
One more word for describing what you lived and felt.
\end{grammarlens}

\subsection*{Your turn}
\begin{itemize}
\raggedright
  \item \emph{Say it:} \emph{paṇipuri}
  \item \emph{Say it:} \emph{paṇipuri}, once more
  \item \emph{Say it:} \emph{eṉ appā muppatu āṇṭukaḷ rayil nilaiyattil paṇipurintār.} — My father worked at the railway station for thirty years.
  \item \emph{Recall:} say the Tamil for an assistant, then the Tamil for a workplace, then say \emph{paṇipuri} again
\end{itemize}

\begin{tcolorbox}[breakable,colback=teal!4,colframe=teal!35!black,title={Before you move on}]
What does \ta{பணிபுரி} mean? (\textquotedblleft{}To be employed; to work\textquotedblright{}.) Say it once more.
\end{tcolorbox}
\section[\texorpdfstring{\ta{செயலாளர்} (ceyalāḷar)}{செயலாளர் (ceyalāḷar)}]{\ta{செயலாளர்} (ceyalāḷar) — a secretary}

\label{lesson:TA-C507-ceyalalar}



\begin{quote}
Before the new one: say the Tamil for a workplace, then the Tamil for to be employed.
\end{quote}

\subsection*{You'll want to know: \ta{செயலாளர்}}
\textbf{\ta{செயலாளர்}} — \emph{ceyalāḷar} — \textquotedblleft{}a secretary\textquotedblright{}.

\begin{quote}
\ta{சங்கத்தின்} \ta{செயலாளர்} \ta{கூட்டத்தை} \ta{நடத்தினார்}.
\end{quote}

\begin{quote}
\emph{caṅkattiṉ ceyalāḷar kūṭṭattai naṭattiṉār.}
\end{quote}

\begin{quote}
The association's secretary ran the meeting.
\end{quote}

Plural: \ta{செயலாளர்கள்}, \emph{ceyalāḷarkaḷ}.

Built like \ta{உதவியாளர்} (\emph{utaviyāḷar}): \ta{செயல்} \textquotedblleft{}work, action\textquotedblright{} + -\ta{ஆளர்} \textquotedblleft{}one who …\textquotedblright{}.

\begin{cousinweb}
\textbf{\ta{செயலாளர்}} (\emph{ceyalāḷar}) is built from native parts: \ta{செயல்} (\emph{ceyal}), \textquotedblleft{}action\textquotedblright{}, from \ta{செய்} (\emph{cey}), \textquotedblleft{}to do\textquotedblright{} (DEDR 1957), and -\ta{ஆளர்}, \textquotedblleft{}one who does it\textquotedblright{}.
\end{cousinweb}

\begin{grammarlens}[title={What you've built}]
One more word for describing what you lived and felt.
\end{grammarlens}

\subsection*{Your turn}
\begin{itemize}
\raggedright
  \item \emph{Say it:} \emph{ceyalāḷar}
  \item \emph{Say it:} \emph{ceyalāḷar}, once more
  \item \emph{Say it:} \emph{caṅkattiṉ ceyalāḷar kūṭṭattai naṭattiṉār.} — The association's secretary ran the meeting.
  \item \emph{Recall:} say the Tamil for a workplace, then the Tamil for to be employed, then say \emph{ceyalāḷar} again
\end{itemize}

\begin{tcolorbox}[breakable,colback=teal!4,colframe=teal!35!black,title={Before you move on}]
What does \ta{செயலாளர்} mean? (\textquotedblleft{}A secretary\textquotedblright{}.) Say it once more.
\end{tcolorbox}
\section[\texorpdfstring{\ta{கோப்பு} (kōppu)}{கோப்பு (kōppu)}]{\ta{கோப்பு} (kōppu) — a file; a folder}

\label{lesson:TA-C507-koppu}



\begin{quote}
Before the new one: say the Tamil for to be employed, then the Tamil for a secretary.
\end{quote}

\subsection*{You'll want to know: \ta{கோப்பு}}
\textbf{\ta{கோப்பு}} — \emph{kōppu} — \textquotedblleft{}a file; a folder\textquotedblright{}.

\begin{quote}
\ta{மேலாளர்} \ta{அந்தக்} \ta{கோப்பை} \ta{என்} \ta{மேசையில்} \ta{வைத்தார்}.
\end{quote}

\begin{quote}
\emph{mēlāḷar antak kōppai eṉ mēcaiyil vaittār.}
\end{quote}

\begin{quote}
The manager put that file on my desk.
\end{quote}

It names a paper folder and a computer file alike: \ta{கோப்பைத்} \ta{திற}, to open a file.

\begin{cousinweb}
\textbf{\ta{கோப்பு}} (\emph{kōppu}) is native Dravidian: DEDR 2176 gives \ta{கோ} (\emph{kō}) and \ta{கோர்} (\emph{kōr}), \textquotedblleft{}to string, to thread\textquotedblright{}, and lists \ta{கோப்பு} for a stringing, a threading. The office sense, a file of papers, is a later use.
\end{cousinweb}

\begin{grammarlens}[title={What you've built}]
One more word for describing what you lived and felt.
\end{grammarlens}

\subsection*{Your turn}
\begin{itemize}
\raggedright
  \item \emph{Say it:} \emph{kōppu}
  \item \emph{Say it:} \emph{kōppu}, once more
  \item \emph{Say it:} \emph{mēlāḷar antak kōppai eṉ mēcaiyil vaittār.} — The manager put that file on my desk.
  \item \emph{Recall:} say the Tamil for to be employed, then the Tamil for a secretary, then say \emph{kōppu} again
\end{itemize}

\begin{tcolorbox}[breakable,colback=teal!4,colframe=teal!35!black,title={Before you move on}]
What does \ta{கோப்பு} mean? (\textquotedblleft{}A file; a folder\textquotedblright{}.) Say it once more.
\end{tcolorbox}
\section[\texorpdfstring{\ta{மேற்பார்வையிடு} (mēṟpārvaiyiṭu)}{மேற்பார்வையிடு (mēṟpārvaiyiṭu)}]{\ta{மேற்பார்வையிடு} (mēṟpārvaiyiṭu) — to supervise}

\label{lesson:TA-C507-merparvaiyitu}



\begin{quote}
Before the new one: say the Tamil for a secretary, then the Tamil for a file.
\end{quote}

\subsection*{You'll want to know: \ta{மேற்பார்வையிடு}}
\textbf{\ta{மேற்பார்வையிடு}} — \emph{mēṟpārvaiyiṭu} — \textquotedblleft{}to supervise\textquotedblright{}.

\begin{quote}
\ta{ஒரு} \ta{மூத்த} \ta{பொறியாளர்} \ta{எங்கள்} \ta{வேலையை} \ta{மேற்பார்வையிடுகிறார்}.
\end{quote}

\begin{quote}
\emph{oru mūtta poṟiyāḷar eṅkaḷ vēlaiyai mēṟpārvaiyiṭukiṟār.}
\end{quote}

\begin{quote}
A senior engineer supervises our work.
\end{quote}

Past (I …): \ta{மேற்பார்வையிட்டேன்}, \emph{mēṟpārvaiyiṭṭēṉ}.

The work watched over takes -\ta{ஐ}: \ta{வேலையை} \ta{மேற்பார்வையிடு}. The noun is \ta{மேற்பார்வை} (\emph{mēṟpārvai}), supervision.

\begin{cousinweb}
\textbf{\ta{மேற்பார்வையிடு}} (\emph{mēṟpārvaiyiṭu}) is a native compound: \ta{மேல்} (\emph{mēl}), \textquotedblleft{}above\textquotedblright{}, \ta{பார்வை} (\emph{pārvai}), \textquotedblleft{}a look\textquotedblright{}, from \ta{பார்} (\emph{pār}), \textquotedblleft{}to see\textquotedblright{} (DEDR 4091), and the light verb \ta{இடு}, \textquotedblleft{}to put\textquotedblright{}.
\end{cousinweb}

\begin{grammarlens}[title={What you've built}]
One more word for describing what you lived and felt.
\end{grammarlens}

\subsection*{Your turn}
\begin{itemize}
\raggedright
  \item \emph{Say it:} \emph{mēṟpārvaiyiṭu}
  \item \emph{Say it:} \emph{mēṟpārvaiyiṭu}, once more
  \item \emph{Say it:} \emph{oru mūtta poṟiyāḷar eṅkaḷ vēlaiyai mēṟpārvaiyiṭukiṟār.} — A senior engineer supervises our work.
  \item \emph{Recall:} say the Tamil for a secretary, then the Tamil for a file, then say \emph{mēṟpārvaiyiṭu} again
\end{itemize}

\begin{tcolorbox}[breakable,colback=teal!4,colframe=teal!35!black,title={Before you move on}]
What does \ta{மேற்பார்வையிடு} mean? (\textquotedblleft{}To supervise\textquotedblright{}.) Say it once more.
\end{tcolorbox}
